\scanpage{11}
% Typed face read after a 180-degree rotation.
\begin{center}- 137 -\end{center}
\hfill n° 391

injectif, et il résulte alors d'\uline{Alg.} \uline{comm.}, chap. III, § 2, n° 8, cor. 1 du th. 1, que $u$ est lui-même \uline{injectif}, ce qui achève de prouver b).

Le fait que b) entraîne c) résulte aussitôt de ce que $k=S'/\underline m'$ est un $S'$-module de dimension projective $r$, donc $\operatorname{Tor}^{S'}_j(k,M)=0$ pour tout $S'$-module $M$ et pour tout $j>r$.

Enfin, on sait que pour tout idéal premier $\underline p$ de $S$, on a $(\operatorname{Tor}^S_j(M,N))_{\underline p}=\operatorname{Tor}^{S_{\underline p}}_j(M_{\underline p},N_{\underline p})$ quels que soient les $S$-modules $M$, $N$. Appliquant ceci à $M=N=k$ et $\underline p=\underline m$, et notant que $k$ est un $S_{\underline m}$-module, donc aussi $\operatorname{Tor}^S_j(k,k)$, on voit que l'on a $\operatorname{Tor}^S_j(k,k)=\operatorname{Tor}^{S_{\underline m}}_j(k,k)$. Le fait que c) entraîne a) résulte alors de la théorie des anneaux locaux réguliers, (\uline{Alg.} \uline{comm.}, chap. VIII).

COROLLAIRE.-- \uline{Si les conditions de la prop.} 1 \uline{sont vérifiées}, $S$ \uline{est un anneau noethérien régulier de dimension} $r$ ; \uline{pour toute famille d'éléments homogènes} $P_i$ \uline{de} $S$ \uline{vérifiant la condition de la prop.} 1 b), \uline{les degrés} $p_i=\deg(P_i)$ \uline{sont les mêmes à l'ordre près}, \uline{et la série de Poincaré de la} $k$-\uline{algèbre graduée} $S$, (\uline{Alg.}, chap. IV, 3e éd.), \uline{est donnée par} :
\[
(1)\qquad P_S(T)=\prod_{i=1}^r(1-T^{p_i})^{-1}.
\]

La première assertion est vraie pour $S'$, (\uline{Alg.} \uline{comm.}, chap. VIII), donc pour l'algèbre $S$ qui lui est isomorphe. En outre les monômes $P_1^{\nu_1}\ldots P_r^{\nu_r}$ sont linéairement indépendants sur $k$, donc $\dim_k(S_n)$ est égal au nombre de familles d'entiers $(\nu_i)_{1\leq i\leq r}$ telles que $\sum_{i=1}^r\nu_i p_i=n$, ou encore au coefficient de $T^n$ au second membre de (1), ce qui démontre cette formule par définition de la série de Poincaré. Ceci prouve déjà l'invariance des degrés $p_i$ ; on peut aussi le voir directement en notant que les $P_i$ appartiennent nécessairement à $\underline m-\underline m^2$, (sans
\mspage{10}
